% ===== PRÉAMBULE CANONIQUE QCM — PHYSIQUE =====
% Sert à compiler controle.tex ET à la revalidation automatique du JSON
% (valider_qcm.py). NE PAS MODIFIER : packages figés.
% La bibliothèque TikZ "babel" est indispensable (conflit babel-french/circuitikz).
\documentclass[11pt,a4paper]{article}
\usepackage[utf8]{inputenc}\usepackage[T1]{fontenc}\usepackage{lmodern}
\usepackage[french]{babel}
\usepackage{amsmath,amssymb,mathtools}\usepackage{icomma}
\usepackage{siunitx}\sisetup{locale=FR}
\usepackage{xcolor}\usepackage{colortbl}
\usepackage{tikz}\usepackage{pgfplots}\pgfplotsset{compat=1.18}
\usepackage[siunitx,europeanresistors]{circuitikz}
\usepackage{enumitem}\usepackage{array,booktabs}
\usepackage[a4paper,margin=2.2cm]{geometry}
\usetikzlibrary{arrows.meta,calc,angles,quotes,positioning,patterns,%
  backgrounds,shapes.geometric,decorations.pathmorphing,%
  decorations.markings,intersections,babel}
\newcommand{\vv}[1]{\vec{#1}}\newcommand{\ds}{\displaystyle}
\newcommand{\R}{\mathbb{R}}\newcommand{\N}{\mathbb{N}}\newcommand{\Z}{\mathbb{Z}}
% ==============================================
\begin{document}

\begin{center}
{\Large\bfseries QCM type concours --- Physique}\\[2pt]
{\large Fiche 08 --- Lentilles et instruments d'optique simples}\\[2pt]
{\itshape 10 questions, niveau concours difficile --- une seule réponse correcte par question.}
\end{center}
\vspace{4pt}
{\footnotesize\itshape Conventions de la fiche : $\dfrac1f=\dfrac1u+\dfrac1v$ ;\ $\gamma=\dfrac{i}{o}=-\dfrac{v}{u}$ ;\
$u>0$ objet réel, $v>0$ image réelle / $v<0$ virtuelle, $f>0$ convergente / $f<0$ divergente,
$i>0$ droite / $i<0$ renversée ; vergence $C=1/f$ ($f$ en mètres, en \si{\per\metre}=dioptrie).}
\vspace{6pt}\hrule\vspace{10pt}

% ---------------------------------------------------------------------
\noindent\textbf{PHYS-F08-Q01 --- Relation de conjugaison : déterminer $f$}\par
Un objet réel lumineux est placé à $u=\qty{24}{\centi\metre}$ devant une lentille mince convergente.
Une image nette est recueillie sur un écran situé à $v=\qty{8}{\centi\metre}$ derrière la lentille
(schéma sans échelle). Quelle est la distance focale $f$ de la lentille ?

\begin{center}
\begin{tikzpicture}[>=Stealth,line cap=round,scale=0.9]
  \draw[gray!60,dashed] (-4.2,0) -- (3.2,0);
  \draw[blue!70,line width=1.3pt,<->] (0,-1.7) -- (0,1.7);
  \node[blue!70,above] at (0,1.7) {\footnotesize L};
  \fill (0,0) circle (1.1pt) node[below right] {\footnotesize $O$};
  \draw[red!75,line width=1.4pt,->] (-3,0) -- (-3,1.1);
  \node[red!75,above] at (-3,1.1) {\footnotesize $B$};
  \node[red!75,below] at (-3,0) {\footnotesize $A$};
  \draw[black,line width=2pt] (1,-1.3) -- (1,1.3);
  \node[below] at (1,-1.3) {\footnotesize écran};
  \draw[<->,gray] (-3,-2.05) -- (0,-2.05);
  \node[gray,below] at (-1.5,-2.05) {\footnotesize $u=24$ cm};
  \draw[<->,gray] (0,-2.05) -- (1,-2.05);
  \node[gray,below] at (0.5,-2.05) {\footnotesize $v=8$ cm};
\end{tikzpicture}
\end{center}

\begin{enumerate}[label=\Alph*)]
\item $\qty{8}{\centi\metre}$
\item $\qty{12}{\centi\metre}$
\item $\qty{6}{\centi\metre}$
\item $\qty{16}{\centi\metre}$
\end{enumerate}
\textit{Bonne réponse : C}\par
Objet et image réels : $\dfrac1f=\dfrac1u+\dfrac1v=\dfrac1{24}+\dfrac18=\dfrac{1}{24}+\dfrac{3}{24}
=\dfrac{4}{24}=\dfrac16$, donc $f=\qty{6}{\centi\metre}$. Cohérent : $u=\qty{24}{\centi\metre}>2f=\qty{12}{\centi\metre}$,
l'image est réelle et réduite ($\gamma=-v/u=-1/3$), régime de l'objectif photographique.\par
\textit{Pièges :}
\begin{itemize}
\item A: a négligé le terme $1/u$ (objet supposé lointain), d'où $f\approx v=\qty{8}{\centi\metre}$.
\item B: a soustrait au lieu d'additionner, $\dfrac1f=\dfrac1v-\dfrac1u=\dfrac{1}{12}$, soit $\qty{12}{\centi\metre}$ ($=2f$).
\item D: a pris la moyenne $(u+v)/2=\qty{16}{\centi\metre}$.
\end{itemize}
\vspace{8pt}\hrule\vspace{8pt}

% ---------------------------------------------------------------------
\noindent\textbf{PHYS-F08-Q02 --- Nature de l'image par la position de l'objet}\par
Une lentille convergente donne, d'un objet réel, une image \textbf{réelle et plus grande} que l'objet.
Dans quelle zone se situe l'objet ?
\begin{enumerate}[label=\Alph*)]
\item au-delà de $2F$ (soit $u>2f$)
\item entre $F$ et $2F$ (soit $f<u<2f$)
\item entre le centre optique $O$ et $F$ (soit $u<f$)
\item exactement au foyer objet $F$ (soit $u=f$)
\end{enumerate}
\textit{Bonne réponse : B}\par
Une image \textbf{réelle} d'un objet réel impose $u>f$ (l'objet est au-delà du foyer). Une image
\textbf{agrandie} ($|\gamma|>1$) impose de plus $u<2f$. Les deux conditions donnent $f<u<2f$ : l'image
est alors réelle, renversée, agrandie et située au-delà de $2F'$ (régime du projecteur).\par
\textit{Pièges :}
\begin{itemize}
\item A: $u>2f$ donne bien une image réelle mais \textbf{réduite} (objectif photo), pas agrandie.
\item C: $u<f$ donne une image \textbf{virtuelle} agrandie (loupe), donc non réelle.
\item D: $u=f$ envoie l'image à l'infini (les rayons émergents sont parallèles).
\end{itemize}
\vspace{8pt}\hrule\vspace{8pt}

% ---------------------------------------------------------------------
\noindent\textbf{PHYS-F08-Q03 --- Lentille divergente et grandissement}\par
Une lentille divergente de distance focale $f=\qty{-20}{\centi\metre}$ donne, d'un objet réel, une image
droite \textbf{trois fois plus petite} que l'objet. À quelle distance $u$ de la lentille l'objet est-il placé ?
\begin{enumerate}[label=\Alph*)]
\item $\qty{40}{\centi\metre}$
\item $\qty{60}{\centi\metre}$
\item $\qty{20}{\centi\metre}$
\item $\qty{80}{\centi\metre}$
\end{enumerate}
\textit{Bonne réponse : A}\par
Une divergente donne toujours une image virtuelle, droite et réduite : ici $\gamma=+\tfrac13$.
Or $\gamma=-\dfrac{v}{u}=+\dfrac13\Rightarrow v=-\dfrac{u}{3}$. La conjugaison donne
$\dfrac1f=\dfrac1u+\dfrac1v=\dfrac1u-\dfrac3u=-\dfrac2u$, d'où $u=-2f=-2\times(-20)=\qty{40}{\centi\metre}$.\par
\textit{Pièges :}
\begin{itemize}
\item B: a écrit directement $u=3|f|=\qty{60}{\centi\metre}$ (confusion « trois fois plus petite » $\to$ « objet à $3f$ »).
\item C: a placé l'objet au foyer, $u=|f|=\qty{20}{\centi\metre}$.
\item D: a oublié que $v<0$ (a pris $v=+u/3$), d'où $u=4|f|=\qty{80}{\centi\metre}$.
\end{itemize}
\vspace{8pt}\hrule\vspace{8pt}

% ---------------------------------------------------------------------
\noindent\textbf{PHYS-F08-Q04 --- Vergence et piège des unités}\par
Quelle est la vergence d'une lentille convergente de distance focale $f=\qty{20}{\centi\metre}$ ?
\begin{enumerate}[label=\Alph*)]
\item $\qty{0.05}{\per\metre}$
\item $\qty{0.2}{\per\metre}$
\item $\qty{50}{\per\metre}$
\item $\qty{5}{\per\metre}$
\end{enumerate}
\textit{Bonne réponse : D}\par
La vergence $C=1/f$ n'est valable que si $f$ est exprimée \textbf{en mètres} : $f=\qty{0.20}{\metre}$, donc
$C=\dfrac{1}{0{,}20}=\qty{5}{\per\metre}$ (soit $+5$ dioptries), lentille convergente.\par
\textit{Pièges :}
\begin{itemize}
\item A: a oublié la conversion en mètres, $1/20=\qty{0.05}{\per\metre}$.
\item B: a donné $f$ (en mètres) au lieu de $1/f$.
\item C: conversion erronée, $f$ pris à $\qty{0.020}{\metre}$ (comme si \qty{20}{\centi\metre} valait \qty{20}{\milli\metre}), $1/0{,}02=\qty{50}{\per\metre}$.
\end{itemize}
\vspace{8pt}\hrule\vspace{8pt}

% ---------------------------------------------------------------------
\noindent\textbf{PHYS-F08-Q05 --- Lentilles minces accolées}\par
On accole (centres optiques confondus) une lentille convergente de vergence $C_1=+\qty{8}{\per\metre}$ et une
lentille divergente de distance focale $f_2=\qty{-50}{\centi\metre}$. Quelle est la distance focale
$f_{\text{éq}}$ du système équivalent ?
\begin{enumerate}[label=\Alph*)]
\item $\qty{-37.5}{\centi\metre}$
\item $\qty{12.5}{\centi\metre}$
\item $\qty{16.7}{\centi\metre}$
\item $\qty{10}{\centi\metre}$
\end{enumerate}
\textit{Bonne réponse : C}\par
Pour des lentilles accolées les vergences s'ajoutent. $C_2=\dfrac{1}{f_2}=\dfrac{1}{-0{,}50}=\qty{-2}{\per\metre}$,
donc $C_{\text{éq}}=C_1+C_2=8-2=\qty{6}{\per\metre}$, et $f_{\text{éq}}=\dfrac16\,\text{m}\approx\qty{16.7}{\centi\metre}$
(système convergent).\par
\textit{Pièges :}
\begin{itemize}
\item A: a additionné les distances focales $f_1+f_2$ (avec $f_1=1/8=\qty{12.5}{\centi\metre}$), $12{,}5-50=\qty{-37.5}{\centi\metre}$.
\item B: a oublié de convertir $f_2$ (\qty{-50}{\centi\metre} traité comme $C_2\approx0$), d'où $f_{\text{éq}}\approx f_1=\qty{12.5}{\centi\metre}$.
\item D: a oublié le signe de la divergente ($C_2=+2$), $C_{\text{éq}}=\qty{10}{\per\metre}\Rightarrow f=\qty{10}{\centi\metre}$.
\end{itemize}
\vspace{8pt}\hrule\vspace{8pt}

% ---------------------------------------------------------------------
\noindent\textbf{PHYS-F08-Q06 --- L'\oe il : vergence en accommodation}\par
L'\oe il normal est modélisé par une lentille convergente qui forme l'image sur la rétine, située à
$\qty{17}{\milli\metre}$ du centre optique. L'\oe il accommode sur un objet placé à $\qty{17}{\centi\metre}$.
En appliquant $\dfrac1f=\dfrac1u+\dfrac1v$, quelle est la vergence $C=1/f$ de l'\oe il dans cet état ?
\begin{enumerate}[label=\Alph*)]
\item $\qty{58.8}{\per\metre}$
\item $\qty{52.9}{\per\metre}$
\item $\qty{5.9}{\per\metre}$
\item $\qty{64.7}{\per\metre}$
\end{enumerate}
\textit{Bonne réponse : D}\par
L'image se forme sur la rétine ($v=\qty{0.017}{\metre}$, réelle) pour un objet à $u=\qty{0.17}{\metre}$ :
$C=\dfrac{1}{u}+\dfrac{1}{v}=\dfrac{1}{0{,}17}+\dfrac{1}{0{,}017}=5{,}88+58{,}82\approx\qty{64.7}{\per\metre}$
(64,7 dioptries).\par
\textit{Pièges :}
\begin{itemize}
\item A: n'a gardé que le terme image $1/v=1/0{,}017\approx\qty{58.8}{\per\metre}$ (oubli de l'objet).
\item B: a soustrait, $1/v-1/u=58{,}82-5{,}88\approx\qty{52.9}{\per\metre}$ (au lieu d'additionner).
\item C: n'a gardé que le terme objet $1/u=1/0{,}17\approx\qty{5.9}{\per\metre}$ (oubli de l'image).
\end{itemize}
\vspace{8pt}\hrule\vspace{8pt}

% ---------------------------------------------------------------------
\noindent\textbf{PHYS-F08-Q07 --- Correction de la myopie}\par
Une personne myope ne voit pas nettement au-delà de $\qty{90}{\centi\metre}$ (punctum remotum). Pour voir
nettement un objet situé \textbf{à l'infini}, on place au contact de l'\oe il un verre correcteur. Quelle est
sa distance focale ?
\begin{enumerate}[label=\Alph*)]
\item $\qty{-90}{\centi\metre}$
\item $\qty{90}{\centi\metre}$
\item $\qty{-45}{\centi\metre}$
\item $\qty{-180}{\centi\metre}$
\end{enumerate}
\textit{Bonne réponse : A}\par
Le verre doit donner de l'objet à l'infini une image virtuelle située au punctum remotum, à
$v=\qty{-0.90}{\metre}$. Avec $u\to\infty$ : $\dfrac1f=0+\dfrac{1}{-0{,}90}\Rightarrow f=\qty{-90}{\centi\metre}$ :
lentille \textbf{divergente} ($C\approx-1{,}1$ dioptrie). Un myope porte des verres divergents.\par
\textit{Pièges :}
\begin{itemize}
\item B: a choisi une lentille \textbf{convergente} ($+90$), qui corrige l'hypermétrope, pas le myope.
\item C: a placé l'image à la moitié du punctum remotum, $\qty{-45}{\centi\metre}$.
\item D: a doublé le punctum remotum, $\qty{-180}{\centi\metre}$.
\end{itemize}
\vspace{8pt}\hrule\vspace{8pt}

% ---------------------------------------------------------------------
\noindent\textbf{PHYS-F08-Q08 --- La loupe : position de l'image}\par
Une loupe (lentille convergente) a une distance focale $f=\qty{5}{\centi\metre}$. On place un petit objet à
$u=\qty{4}{\centi\metre}$ de la loupe (schéma sans échelle). Où se forme l'image, comptée algébriquement
depuis la loupe ?

\begin{center}
\begin{tikzpicture}[>=Stealth,line cap=round,scale=0.9]
  \draw[gray!60,dashed] (-4,0) -- (3.2,0);
  \draw[blue!70,line width=1.3pt,<->] (0,-1.6) -- (0,1.6);
  \fill (0,0) circle (1.1pt) node[below right] {\footnotesize $O$};
  \fill (-2,0) circle (1.3pt) node[below] {\footnotesize $F$};
  \fill (2,0) circle (1.3pt) node[below] {\footnotesize $F'$};
  \draw[red!75,line width=1.4pt,->] (-1.6,0) -- (-1.6,0.9);
  \node[red!75,above] at (-1.6,0.9) {\footnotesize objet};
  \draw[<->,gray] (-2,-1.9) -- (0,-1.9);
  \node[gray,below] at (-1,-1.9) {\footnotesize $f=5$ cm};
\end{tikzpicture}
\end{center}

\begin{enumerate}[label=\Alph*)]
\item $\qty{20}{\centi\metre}$
\item $\qty{-20}{\centi\metre}$
\item $\qty{-2.2}{\centi\metre}$
\item $\qty{-5}{\centi\metre}$
\end{enumerate}
\textit{Bonne réponse : B}\par
$\dfrac1v=\dfrac1f-\dfrac1u=\dfrac15-\dfrac14=\dfrac{4-5}{20}=-\dfrac{1}{20}\Rightarrow v=\qty{-20}{\centi\metre}$.
Le signe négatif confirme une image \textbf{virtuelle}, à \qty{20}{\centi\metre} devant la loupe, du même côté
que l'objet ; elle est droite et agrandie ($\gamma=-v/u=+5$).\par
\textit{Pièges :}
\begin{itemize}
\item A: erreur de signe, image déclarée réelle ($v=+\qty{20}{\centi\metre}$) alors qu'une loupe donne une image virtuelle.
\item C: a additionné, $\dfrac{1}{|v|}=\dfrac1f+\dfrac1u=\dfrac{9}{20}$, d'où $\qty{-2.2}{\centi\metre}$.
\item D: a pris $|v|=f$ (image supposée au foyer), $\qty{-5}{\centi\metre}$.
\end{itemize}
\vspace{8pt}\hrule\vspace{8pt}

% ---------------------------------------------------------------------
\noindent\textbf{PHYS-F08-Q09 --- Rayon particulier sur une lentille divergente}\par
Un rayon incident, parallèle à l'axe optique, frappe une lentille \textbf{divergente} (schéma). Quelle est la
trajectoire du rayon émergent ?

\begin{center}
\begin{tikzpicture}[>=Stealth,line cap=round,scale=0.9]
  \draw[gray!60,dashed] (-4.2,0) -- (3.2,0);
  \draw[blue!70,line width=1.3pt,>-<] (0,-1.6) -- (0,1.6);
  \fill (0,0) circle (1.1pt) node[below right] {\footnotesize $O$};
  \fill (-1.6,0) circle (1.3pt) node[below] {\footnotesize $F'$};
  \fill (1.6,0) circle (1.3pt) node[below] {\footnotesize $F$};
  \draw[red!75,line width=1pt,->] (-4,0.9) -- (0,0.9);
  \node[red!75,above] at (-3,0.9) {\footnotesize rayon incident};
  \node[blue!70,below] at (0,-1.9) {\footnotesize lentille divergente};
\end{tikzpicture}
\end{center}

\begin{enumerate}[label=\Alph*)]
\item Il émerge en passant par le foyer image $F'$ situé de l'autre côté de la lentille (après elle).
\item Il traverse la lentille sans être dévié et poursuit parallèlement à l'axe.
\item Il émerge en s'écartant de l'axe, son prolongement passant par le foyer image $F'$ situé du côté de l'objet (avant la lentille).
\item Il émerge en passant par le foyer objet $F$.
\end{enumerate}
\textit{Bonne réponse : C}\par
Pour une lentille divergente, un faisceau incident parallèle à l'axe émerge en s'écartant : les rayons
émergents \emph{semblent provenir} du foyer image $F'$, lequel est situé \textbf{du côté de l'objet}
(à gauche de $O$). C'est donc le \emph{prolongement} du rayon émergent qui passe par $F'$.\par
\textit{Pièges :}
\begin{itemize}
\item A: décrit le comportement d'une \textbf{convergente} ($F'$ réel, après la lentille), pas d'une divergente.
\item B: seul le rayon passant par le centre optique $O$ n'est pas dévié ; un rayon parallèle à l'axe est nécessairement dévié.
\item D: confond foyer image et foyer objet ; c'est un rayon passant \emph{par} $F$ qui ressort parallèle (situation inverse).
\end{itemize}
\vspace{8pt}\hrule\vspace{8pt}

% ---------------------------------------------------------------------
\noindent\textbf{PHYS-F08-Q10 --- Système afocal de deux lentilles}\par
Un système afocal est formé de deux lentilles convergentes de distances focales $f_1=\qty{20}{\centi\metre}$ et
$f_2=\qty{5}{\centi\metre}$, alignées sur le même axe, de sorte qu'un faisceau incident parallèle à l'axe
ressorte parallèle (image finale à l'infini). Quelle doit être la distance $d$ entre les deux lentilles ?

\begin{center}
\begin{tikzpicture}[>=Stealth,line cap=round,scale=0.85]
  \draw[gray!60,dashed] (-1.3,0) -- (6.3,0);
  \draw[blue!70,line width=1.3pt,<->] (0,-1.5) -- (0,1.5);
  \node[blue!70,above] at (0,1.5) {\footnotesize $L_1$};
  \draw[blue!70,line width=1.3pt,<->] (4,-1.2) -- (4,1.2);
  \node[blue!70,above] at (4,1.2) {\footnotesize $L_2$};
  \draw[red!70,->] (-1.3,0.8) -- (0,0.8);
  \draw[red!70] (0,0.8) -- (2.4,0);
  \fill (2.4,0) circle (1.1pt);
  \node[below] at (2.4,-0.05) {\footnotesize $F'_1=F_2$};
  \draw[red!70,->] (2.4,0) -- (4,0.55);
  \draw[red!70,->] (4,0.55) -- (6.1,0.55);
  \draw[<->,gray] (0,-1.75) -- (4,-1.75);
  \node[gray,below] at (2,-1.75) {\footnotesize $d = ?$};
\end{tikzpicture}
\end{center}

\begin{enumerate}[label=\Alph*)]
\item $\qty{25}{\centi\metre}$
\item $\qty{15}{\centi\metre}$
\item $\qty{4}{\centi\metre}$
\item $\qty{100}{\centi\metre}$
\end{enumerate}
\textit{Bonne réponse : A}\par
Le faisceau parallèle converge au foyer image $F'_1$ de la première lentille, à $f_1$ derrière $L_1$. Pour
ressortir parallèle, ce point doit coïncider avec le foyer objet $F_2$ de la seconde, à $f_2$ devant $L_2$.
Les foyers sont donc confondus : $d=f_1+f_2=20+5=\qty{25}{\centi\metre}$.\par
\textit{Pièges :}
\begin{itemize}
\item B: a soustrait les focales, $f_1-f_2=\qty{15}{\centi\metre}$.
\item C: a appliqué la formule des lentilles \textbf{accolées} $\dfrac{f_1 f_2}{f_1+f_2}=\qty{4}{\centi\metre}$ (valable seulement centres confondus, $d=0$).
\item D: a multiplié les focales, $f_1\times f_2=\qty{100}{\centi\metre}$ (sans diviser).
\end{itemize}

\end{document}
